% ---------------------------------------------------------------------------
% Atom localisation from the multislice-ptychographic object volume
% Requires: amsmath (\eqref, \lVert). Citations use natbib (\citep) or biblatex;
% references in atomfind_refs.bib
% ---------------------------------------------------------------------------
\subsection{Atom localisation with a measured three-dimensional point-spread function}
\label{sec:atomfind}

\paragraph{Forward model.}
We treat the reconstructed multislice object as a sparse superposition of identical
single-atom responses,
\begin{equation}
  \phi(\mathbf{r}) \;=\; \sum_{i=1}^{N} \beta_i\, K(\mathbf{r}-\mathbf{r}_i)
  \;+\; b(\mathbf{r}) \;+\; \varepsilon(\mathbf{r}),
  \qquad \beta_i \ge 0,
  \label{eq:forward}
\end{equation}
with $\phi$ the per-layer-median-subtracted phase, $(\beta_i,\mathbf{r}_i)$ the amplitude
and three-dimensional position of atom $i$, $b$ a slowly varying background removed by a
high-pass filter, and $\varepsilon$ the residual. The
reconstruction is a nonlinear inverse problem, so Eq.~\eqref{eq:forward} is an empirical
linearisation rather than a derived identity; we test it directly by fixing
$\{\mathbf{r}_i\}$ at the known atomic positions of a simulated structure and solving the
resulting linear least-squares problem for $\{\beta_i\}$, which reproduces measured column
profiles with a correlation of $0.94$.

The kernel $K$ is measured rather than parameterised: a single isolated atom is propagated
through the identical simulation and reconstruction pipeline, and its reconstructed
response is taken as $K$. This is material because $K$ is strongly anisotropic --- full
width at half maximum (FWHM) $\approx 0.1$\,\AA{} in-plane against $2$--$3$\,\AA{} along
the beam, the latter set by the missing-cone transfer of single-projection ptychography
\citep{Chen2021} --- and because its axial width exceeds the projected interatomic spacing
($1.9$--$3.9$\,\AA{} here). Neighbouring responses therefore overlap, the design is not
separable, and localisation cannot be reduced to independent peak fits.

Estimation proceeds in two stages: greedy selection of the support and model order $N$,
then continuous refinement of $\{\beta_i,\mathbf{r}_i\}$ against
Eq.~\eqref{eq:forward}. Greedy selection followed by local fitting is the established
route to atom localisation in three-dimensional electron microscopy \citep{Xu2015}.

\paragraph{Support selection.}
$K$ is normalised so that $\lVert K\rVert_2 = 1$. For a single source at position
$\mathbf{r}$, the least-squares amplitude is then $\hat\beta(\mathbf{r}) = \langle
R, K(\cdot-\mathbf{r})\rangle = (K\star R)(\mathbf{r})$, the matched filter applied to the
current residual $R$; and the position maximising $\hat\beta$ is the position maximising
the achievable reduction in residual sum of squares. Selection is therefore greedy forward
selection on the dictionary $\{K(\cdot-\mathbf{r})\}$, i.e.\ matching pursuit
\citep{MallatZhang1993}: take $\hat{\mathbf{r}} = \arg\max (K\star R)$, record
$(\hat\beta,\hat{\mathbf{r}})$, update $R \leftarrow R - \hat\beta K(\cdot-\hat{\mathbf{r}})$,
exclude a minimum-separation neighbourhood of $\hat{\mathbf{r}}$ from subsequent selection,
and repeat. The exclusion enforces identifiability: it prevents two dictionary atoms
separated by much less than the width of $K$, whose columns are near-collinear, from both
entering the support.

Iteration stops when $\max(K\star R) < k\,\sigma_{\mathrm{MF}}$, where
$\sigma_{\mathrm{MF}}$ is a median-absolute-deviation estimate of the matched-filter noise
computed over the sub-median voxels of the same subvolume, and $k=2$ throughout. The
threshold is noise-relative rather than absolute because $(K\star R)$ scales with the phase
amplitude of the reconstruction, which differs by a factor of about two between the
reconstructions used here; an absolute threshold does not transfer between them. Because
overlapping responses bias the greedy amplitudes --- each is estimated as though its source
were isolated --- the stopping rule is set for completeness of the support, and the
amplitudes themselves are re-estimated in the next stage.

\paragraph{Refinement.}
Writing $\theta_i = (\beta_i,\mathbf{r}_i)$, we minimise
$\lVert \phi - \sum_i \beta_i K(\cdot-\mathbf{r}_i)\rVert^2$ by block-coordinate descent:
each atom is updated in turn by Gauss--Newton steps on the local residual with the current
models of all other atoms subtracted, so that overlap between neighbours is accounted for
to first order at every update. Derivatives of $K$ with respect to $\mathbf{r}$ are obtained
by finite differences on the measured kernel. Two sweeps are used; the update is confined to
a neighbourhood of each atom, so the cost is linear in $N$. This is model-based fitting of a
peak superposition \citep{DeBacker2016} carried into three dimensions with a measured rather
than parametric $K$. A candidate is retained only if the normalised correlation between the
local data (neighbours removed) and its fitted response exceeds a fixed threshold --- a
goodness-of-fit criterion on shape, independent of amplitude, so that weak but well-formed
responses are not discarded with the poorly formed ones. Estimating depth by fitting the
axial response is established for isolated dopants
\citep{Chen2021,Ishikawa2014}; here it is applied to every atom of a dense lattice.

\paragraph{Species assignment and lattice-constrained detection.}
Columns are typed from their own fitted-amplitude statistics, which are well separated in
this material; within mixed columns the two sublattices are separated by amplitude, with
ambiguous cases resolved by consistency with the local periodicity of the already-assigned
neighbours. Sites predicted by that periodicity --- from the fitted atoms of the same
column, not from the reference structure --- but left empty by the selection stage are
re-examined at a reduced detection threshold with bounded positional and amplitude
tolerances. This deliberately trades false-negative rate against false-positive rate for
the weakest species only, and such detections are flagged in all outputs so that the
constrained and unconstrained populations can be reported separately.

\paragraph{Uncertainty.}
From the final Gauss--Newton step of each atom we form the block covariance
$\hat{C}_i = s_i^2 (\mathbf{J}_i^{\!\top}\mathbf{J}_i)^{-1}$, with $\mathbf{J}_i$ the
Jacobian of Eq.~\eqref{eq:forward} with respect to $\theta_i$ and $s_i^2$ the residual
variance on the local patch. Two properties of $\hat{C}_i$ must be stated. First, it is
block-diagonal by construction: it is evaluated with neighbouring atoms held fixed, so
correlations between the parameters of overlapping atoms are not propagated, and $\hat{C}_i$
is optimistic wherever two responses overlap strongly. Second, on near-noiseless data
$s_i^2$ is dominated by model mismatch --- the difference between the measured $K$ and the
true in-crystal response --- rather than by counting statistics, so $\hat{C}_i$ alone
underestimates the error.

We therefore report a variance decomposition
$\sigma^2 = \sigma_{\mathrm{fit}}^2 + \sigma_{\mathrm{sys}}^2$, where
$\sigma_{\mathrm{fit}}^2 = \mathrm{diag}(\hat{C}_i)$ and $\sigma_{\mathrm{sys}}$ is a
per-species systematic term. $\sigma_{\mathrm{sys}}$ is fixed by empirical coverage: it is
chosen so that $68\%$ of atoms in a simulation with known ground truth satisfy
$\lvert\Delta\rvert\le\sigma$ on each axis, and coverage is then audited separately by
species, axis, depth and detection mode rather than assumed uniform. Because
$\sigma_{\mathrm{sys}}$ represents a localisation limit set by the width of the response, it
is rescaled for a new reconstruction by the measured FWHM of its kernel; empirical coverage
is reported with the coordinates, so a dataset on which the calibration fails to transfer is
identified rather than silently trusted. The resulting precision, $\sigma \approx
\mathrm{FWHM}/10$ per axis, is consistent with the width-over-signal-to-noise scaling
expected for the localisation of a known template and lies well below the resolution.

\paragraph{Comparison with deconvolution.}
The established alternative --- three-dimensional Richardson--Lucy or maximum-entropy
deconvolution of the volume followed by peak detection \citep{Ishizuka2021} --- was run
with the same measured $K$ on the same volume. Both routines sharpen the strongly
scattering columns but recover at most $59\%$ of the oxygen sublattice at their
best-effort thresholds, below the $71\%$ obtained by peak-picking the undeconvolved volume
and well below the $96\%$ (in-bulk) of the model fit. Deconvolution redistributes dynamic
range towards the strongly scattering species and cannot restore the missing-cone
frequencies that separate an oxygen atom from the axial flank of its titanium neighbour;
the reported advantage of these routines is obtained for sparse, isolated features rather
than for a dense sublattice.
% ---------------------------------------------------------------------------
